\scanpage{89}
Un homomorphisme $u:A\to B$ de V.A. est connu quand on connaît la classe de cohomologie associée : son graphe dans $A\times B$, \deleted{\unsure{car cela|car cette|en cette} \illegible{1}} (i.e. si on le connaît en tant que \unsure{classe|dans} de corr. cohomologique $l$-adique), car alors on connaît le graphe $\Gamma\inserted{\subset A\times B}$ de $u$ car $T_{Q_\ell}(\Gamma)\subset T_{Q_\ell}(A\times B)$ est alors connu.

\uline{D} $=$, une ss-V.A. de $A$ est connue quand on connaît la classe de cohomologie.

L'\textquotedblleft interprétation\textquotedblright\ cohomologique de l'isom
\[
 C^1(A\times B)/\deleted{\unsure{C^1}} C^1(A)\times C^1(B)\simeq\operatorname{Hom}^{\unsure{\prime}}(A,B^*)
\]
est bien claire, via
\[
\begin{aligned}
&\textstyle\bigwedge^2(H^1(A)\times H^1(B))(1)/\bigwedge^2 H^1(A)(1)\times\bigwedge^2 H^1(B)(1)
 \simeq H^1(A)\otimes H^1(B)(1)\\[-2pt]
&\hfill \rotatebox{90}{$\unsure{\simeq|=|\rotatebox{-90}{$\wr$}}$}\qquad\\[-2pt]
&\hfill \operatorname{Hom}(T_\ell(A),T_\ell(B^*))
\end{aligned}
\]

\uline{Question} Prouver que pour qu'un \unsure{élément|élt|élts} de $H^1(A)\otimes H^1(B)(1)\subset H^2(A\times B)(1)$ soit algébrique

% archive sheet number: 88
